\section{Conclusion}

\subsection{Summary}

We evaluated the hypothesis that better pre-training data curation produces more generalizable language models, operationalized as a higher MMLU/HellaSwag generalization balance ratio for OLMo-7B (Dolma) over Pythia-6.9B (The Pile) at matched training scale ($\sim$300B tokens). The hypothesis is not supported: Pythia-6.9B achieves the higher ratio (0.565 vs.\ 0.538; difference $= -0.0265$; 95\% CI $[-0.045, -0.007]$; Cohen's $d = -2.732$). The most structurally informative finding is HellaSwag denominator saturation: both models achieve identical commonsense scores (0.4580 each), collapsing the ratio metric into a MMLU proxy and limiting the interpretability of generalization balance comparisons at this scale. The result is consistent with architecture confounds dominating corpus quality effects at 300B tokens, and motivates architecture-matched or trajectory-based designs for future corpus quality studies.

\subsection{Future Directions}

\textbf{Architecture-matched corpus ablations.} The most direct follow-up is to train OLMo-architecture models on both Dolma and The Pile (or a subset thereof) at matched scale, isolating corpus quality as the sole variable. This requires significant compute but is the minimum design for causal corpus quality attribution.

\textbf{Multi-checkpoint trajectory analysis.} Evaluating both models at 100B, 200B, 300B, and 400B+ token checkpoints would reveal whether corpus quality effects emerge at higher scales, or whether the architecture confound is present throughout training. Both Pythia and OLMo provide public intermediate checkpoints that make this analysis feasible without additional training.

\textbf{Denominator benchmark selection.} Future generalization balance studies should identify benchmarks that remain sensitive to corpus quality variation at the model scales and training durations being compared. Benchmarks with steeper scaling curves --- or subject-level decompositions that preserve variance --- may be better suited as ratio denominators than HellaSwag at the 6--8B, 300B token regime.

\textbf{Contamination-controlled evaluation.} Applying n-gram overlap analysis~\cite{shi2024detecting} between training corpora and evaluation benchmarks prior to comparison would bound the contamination alternative explanation and increase the credibility of corpus quality conclusions.

\subsection{Closing}

The reproducibility of this study --- using public checkpoints, a standardized evaluation harness, and fully documented configurations --- is itself a contribution to the infrastructure for controlled LLM evaluation. We release all evaluation scripts and results to support replication and extension. The field's investment in corpus curation is well-motivated by intuition and prior aggregate evidence; our results suggest the next priority is controlled experimental designs that can isolate and measure the effect size of curation quality at scale.
